\section{Phương pháp và triển khai}
\label{sec:method}

\subsection{Tổng quan pipeline}

Pipeline của đề tài được tổ chức thành các bước rời nhau, mỗi bước nằm trong một
file Python riêng dưới thư mục \texttt{src/}, để dễ kiểm thử và chạy được trên cả
Google Colab lẫn máy cá nhân. Trình tự xử lý như sau:

\begin{enumerate}[leftmargin=1.6em,itemsep=1pt]
  \item Đọc file âm thanh từ UrbanSound8K.
  \item Chuyển tín hiệu về mono.
  \item Resample về 22{,}05\,kHz.
  \item Pad hoặc cắt về thời lượng cố định 4{,}0 giây.
  \item Trích xuất Mel spectrogram rồi lấy log để được log-Mel spectrogram.
  \item Chuẩn hóa đặc trưng.
  \item Tạo \texttt{Dataset} và \texttt{DataLoader} cho train, validation, test.
  \item Huấn luyện mô hình (CNN, ResNet18, hoặc ResNet18 + augmentation).
  \item Đánh giá bằng Accuracy, Macro-F1, classification report và confusion
        matrix; với họ ResNet còn áp dụng test-time augmentation.
\end{enumerate}

Một chi tiết kỹ thuật đáng nói: bước trích xuất log-Mel spectrogram được thực
hiện một lần cho toàn bộ dữ liệu và lưu lại thành một tệp đặc trưng
(\textit{feature cache}). Nhờ đó, các lần huấn luyện sau không phải đọc và biến
đổi lại âm thanh, rút ngắn đáng kể thời gian thí nghiệm.

\subsection{Tiền xử lý dữ liệu âm thanh}

\subsubsection{Chuyển về mono}

Do dữ liệu trộn mono và stereo, mọi file được đưa về mono bằng cách lấy trung
bình các kênh. Bước này chuẩn hóa số chiều kênh của đầu vào và giảm chi phí tính
toán.

\subsubsection{Resample về 22{,}05\,kHz}

EDA cho thấy tần số lấy mẫu trải dài từ 8\,kHz đến 192\,kHz. Toàn bộ dữ liệu được
resample về 22{,}05\,kHz -- một mức phổ biến trong xử lý âm thanh, đủ giữ phần
lớn thông tin hữu ích của âm thanh đô thị (chủ yếu nằm ở dải tần thấp và trung)
trong khi vẫn nhẹ hơn nhiều so với 44{,}1\,kHz.

\subsubsection{Chuẩn hóa độ dài về 4{,}0 giây}

Để mọi mẫu có cùng kích thước, độ dài được cố định ở 4{,}0 giây. Với tần số
22{,}05\,kHz, điều này tương ứng $22050 \times 4{,}0 = 88200$ mẫu. Đoạn ngắn hơn
được pad bằng giá trị 0 ở cuối; đoạn dài hơn được cắt về đúng 88200 mẫu.

\subsubsection{Trích xuất log-Mel spectrogram}

Tín hiệu sau chuẩn hóa được chuyển sang Mel spectrogram với các tham số:
$n_{\text{fft}} = 1024$, $\text{hop\_length} = 512$, $n_{\text{mels}} = 128$.
Biên độ sau đó được đưa về thang decibel với $\text{top\_db} = 80$. Kết quả là
một log-Mel spectrogram một kênh kích thước $128 \times 173$ (tần số $\times$
thời gian), đóng vai trò đặc trưng đầu vào cho mọi mô hình.

\subsubsection{Chuẩn hóa đặc trưng}

Mỗi spectrogram được chuẩn hóa về trung bình 0 và độ lệch chuẩn 1. Bước này thu
hẹp chênh lệch biên độ giữa các mẫu, giúp quá trình tối ưu ổn định và hội tụ
nhanh hơn.

\subsection{Tổ chức dữ liệu và \texttt{DataLoader}}

Dữ liệu được chia dựa trên cấu trúc 10 fold sẵn có. Trong toàn bộ thực nghiệm,
fold~1 được dùng làm tập test, fold~2 làm tập validation, tám fold còn lại làm
tập train. Cách chia này đảm bảo tập test hoàn toàn tách khỏi quá trình huấn
luyện và việc lựa chọn mô hình.

Mỗi mẫu do \texttt{Dataset} trả về là một dictionary gồm: đặc trưng log-Mel, nhãn
số, tên lớp, tên file, fold và đường dẫn. Việc giữ kèm tên file và tên lớp tuy
không cần cho huấn luyện nhưng rất tiện khi gỡ lỗi và phân tích các trường hợp dự
đoán sai. \texttt{DataLoader} trả về ba bộ nạp riêng cho train, validation và
test. Ở nhánh có augmentation, chỉ tập train mới được tăng cường; validation và
test luôn giữ nguyên để kết quả đánh giá phản ánh đúng năng lực tổng quát hóa.

\subsection{Mô hình CNN baseline}

CNN baseline là mô hình cơ sở, đóng vai trò mốc tham chiếu cho các cấu hình phức
tạp hơn. Kiến trúc gồm bốn khối tích chập với số kênh tăng dần $16 \to 32 \to 64
\to 128$. Mỗi khối gồm một tầng tích chập $3\times 3$, batch normalization và
ReLU; ba khối đầu kèm max pooling $2\times 2$, khối cuối dùng adaptive average
pooling để thu về một vector đặc trưng 128 chiều. Bộ phân loại gồm một tầng ẩn
$128 \to 64$ với ReLU và dropout $0{,}3$, rồi một tầng $64 \to 10$.

Lý do chọn CNN làm baseline rất rõ ràng: kiến trúc này phù hợp tự nhiên với dữ
liệu hai chiều dạng spectrogram, dễ huấn luyện và ít nhạy với siêu tham số. Nó
cho biết một mô hình đơn giản nhưng phù hợp có thể đạt tới đâu, từ đó làm chuẩn
để đánh giá liệu ResNet18 có thực sự mang lại lợi ích hay không.

\subsection{Mô hình ResNet18}

ResNet18 được dùng như mô hình sâu hơn, đủ năng lực học các đặc trưng phức tạp mà
CNN baseline khó nắm bắt. Đề tài không dùng trọng số tiền huấn luyện, nhằm đánh
giá trung thực khả năng học của kiến trúc trực tiếp từ UrbanSound8K. Hai điều
chỉnh được thực hiện so với ResNet18 gốc dành cho ảnh:

\begin{itemize}[leftmargin=1.4em,itemsep=2pt]
  \item \textbf{Đầu vào một kênh.} Tầng tích chập đầu được sửa để nhận trực tiếp
        spectrogram một kênh, thay vì ba kênh RGB, tránh phải nhân bản kênh một
        cách giả tạo.
  \item \textbf{Stem nhẹ.} Phần stem được giữ ở dạng tích chập $3\times 3$ stride
        1, không hạ kích thước sớm bằng pooling. Trên spectrogram, các cấu trúc
        năng lượng nhỏ theo thời gian hoặc tần số có thể mang ý nghĩa phân biệt,
        nên cần được giữ lại ở giai đoạn đầu.
\end{itemize}

Mô hình gồm bốn tầng (stage), mỗi tầng có hai khối \texttt{BasicBlock} với kết
nối tắt; dropout dạng không gian được chèn trong khối để hạn chế overfitting. Số
kênh nền (\textit{base channels}) được dùng như một tham số dung lượng: cấu hình
ResNet18 không augmentation dùng base channels 48, còn cấu hình có augmentation
nâng lên 64 để bù lại lượng nhiễu mà augmentation đưa vào.

\subsection{Augmentation trong miền spectrogram}

Nhánh augmentation áp dụng các phép biến đổi sau, \textit{chỉ trên tập train}:

\begin{itemize}[leftmargin=1.4em,itemsep=2pt]
  \item thêm nhiễu Gaussian nhẹ;
  \item dịch theo trục thời gian bằng phép \texttt{roll};
  \item frequency masking -- che một dải tần số;
  \item time masking -- che một dải thời gian.
\end{itemize}

Ngoài bốn phép trên ở mức từng mẫu, MixUp được áp dụng ở mức batch: với một xác
suất nhất định, hai mẫu trong batch được trộn tuyến tính cùng với nhãn của chúng.

Cường độ augmentation đã được tinh chỉnh trong quá trình làm. Cấu hình ban đầu
đặt masking khá rộng (che tới khoảng $17\%$ chiều thời gian, hai dải mỗi loại) và
MixUp với xác suất $0{,}5$. Quan sát kết quả cho thấy mức này làm hỏng đặc trưng
của nhóm âm thanh cơ khí vốn ``đều và lặp'' -- che mất một dải thời gian dài đồng
nghĩa với xóa đi chính tính lặp đặc trưng của chúng. Vì vậy, cấu hình cuối cùng
giảm độ rộng masking xuống còn một dải hẹp mỗi loại và hạ xác suất MixUp xuống
$0{,}3$. Đây là một minh họa cụ thể cho nhận định ở Mục~\ref{sec:litreview}:
augmentation quá mạnh có thể phản tác dụng.

\subsection{Test-time augmentation và ensemble}

\subsubsection{Test-time augmentation (TTA)}

Khi đánh giá họ mô hình ResNet, đề tài áp dụng test-time augmentation. Với mỗi
mẫu test, mô hình dự đoán trên ba phiên bản đầu vào -- bản gốc và hai bản dịch
nhẹ theo trục thời gian -- rồi lấy trung bình logit. TTA không phải huấn luyện
lại, chi phí thấp, và giúp dự đoán ổn định hơn trước những dịch chuyển nhỏ về
thời gian vốn không làm thay đổi nhãn.

\subsubsection{Ensemble hai mô hình}

Cấu hình thứ tư của đề tài là ensemble giữa ResNet18 và ResNet18 + augmentation.
Hai mô hình này được huấn luyện độc lập với cấu hình khác nhau nên có \textit{hồ
sơ lỗi} khác nhau: bản không augmentation mạnh ở nhóm âm thanh cơ khí, bản có
augmentation mạnh hơn ở một số lớp ``trầm'' và ít mẫu. Ensemble được thực hiện
đơn giản bằng cách lấy trung bình xác suất softmax (sau TTA) của hai mô hình với
trọng số bằng nhau. Mục tiêu là khai thác tính bù trừ đó mà không cần huấn luyện
thêm.

\subsection{Chiến lược huấn luyện}

Tất cả mô hình được huấn luyện trên tập train, theo dõi trên tập validation và
chỉ checkpoint tốt nhất (theo Accuracy validation) mới được dùng để đánh giá
test. Một số điểm chung trong recipe huấn luyện: hàm mất mát là cross-entropy có
label smoothing; tối ưu bằng Adam/AdamW~\cite{loshchilov2019}; learning rate được
điều chỉnh tự động (CNN dùng \texttt{ReduceLROnPlateau}, họ ResNet dùng cosine
annealing); và có early stopping để dừng khi validation không còn cải thiện.
ResNet còn dùng gradient clipping và mixed-precision để huấn luyện ổn định, nhanh
hơn. Các siêu tham số cụ thể của từng cấu hình được liệt kê ở
Mục~\ref{sec:expdesign}.

\subsection{Chỉ số đánh giá và công cụ theo dõi}

Đề tài dùng ba nhóm chỉ số. \textbf{Accuracy} -- tỷ lệ mẫu dự đoán đúng -- cho
cái nhìn tổng thể nhưng có thể che giấu sự mất cân bằng giữa các lớp.
\textbf{Macro-F1} -- trung bình F1-score của mười lớp -- bổ khuyết cho điều đó:
vì mỗi lớp đóng góp như nhau bất kể số mẫu, một mô hình chỉ giỏi vài lớp dễ sẽ bị
Macro-F1 ``phạt''. \textbf{Confusion matrix} và \textbf{classification report}
cho phép phân tích định tính: nhìn ra chính xác cặp lớp nào hay bị nhầm và lớp
nào kéo điểm chung xuống.

Toàn bộ quá trình huấn luyện và đánh giá được ghi log bằng Weights \& Biases. Công
cụ này lưu lại loss, accuracy, Macro-F1 theo từng epoch, lưu checkpoint và lịch
sử các lần chạy, nhờ đó việc so sánh giữa các cấu hình trở nên minh bạch và tái
lập được.
